% part: intuitionistic-logic
% chapter: semantics
% section: propositions

\documentclass[../../../include/open-logic-section]{subfiles}

\begin{document}

\olfileid{int}{sem}{pro}

\olsection{વિધાનો}

ક્યારેક !!a{relational model}~$\mModel{M}$માં
!!a{formula}~$!A$ જ્યાં સાચું હોય તે જગતોના
ગણને નામ આપવું ઉપયોગી છે. તેને $!A$ દ્વારા
વ્યાખ્યાયિત \emph{વિધાન} કહીએ છીએ અને
$\Prop{M}{!A}$ લખીએ છીએ.

\begin{defn}
  $V$ વિધેયને આગમનાત્મક રીતે
  $\Prop{M}{\cdot} \colon \Frm
  \to \Pow{W}$ વિધેય સુધી આ પ્રમાણે વિસ્તરીએ:
  \begin{enumerate}
  \item $\Prop{M}{\lfalse} = \emptyset$
  \item $\Prop{M}{p} = V(p)$
  \item $\Prop{M}{\lnot !B} = {}$
    \[
      \Setabs{w \in W}{\text{દરેક $w'$ માટે જે $Rww'$ સંતોષે છે, $w'
          \notin \Prop{M}{!B}$}}.
    \]
  \item $\Prop{M}{!B \land !C} = \Prop{M}{!B} \cap \Prop{M}{!C}$.
  \item $\Prop{M}{!B \lor !C} = \Prop{M}{!B} \cup \Prop{M}{!C}$.
  \item $\Prop{M}{!B \lif !C} = {}$
    \[
      \Setabs{w \in W}{\text{દરેક $w'$ માટે જે $Rww'$ સંતોષે છે, જો $w'
          \in \Prop{M}{!B}$ હોય તો $w' \in \Prop{M}{!C}$}}.
    \]
  \end{enumerate}
  \sourcecorrection{OLMD-185}{મૂળમાં આખા વિધેયની
  પ્રકાર-ઘોષણામાં $\Prop{M}{!A}$ લખ્યું છે, જોકે
  $!A$ અહીં વિધેયનું આર્ગ્યુમેન્ટ છે; તેથી
  $\Prop{M}{\cdot}$ રાખ્યું છે.}
\end{defn}

\begin{prop}
  \ollabel{prop:propositions-true}
  $\mSat{M}{!A}[w]$ ત્યારે અને માત્ર ત્યારે જ જ્યારે
  $w \in \Prop{M}{!A}$.
  \sourcecorrection{OLMD-186}{મૂળ વિધાનમાં
  $\mModel{M}{!A}[w]$ છે; અગાઉ વ્યાખ્યાયિત
  જગત-ખાતે સત્યતાની સંજ્ઞા $\mSat{M}{!A}[w]$
  અહીં જરૂરી છે.}
\end{prop}

\begin{proof}
  સ્વાધ્યાય.
\end{proof}

\begin{prob}
  \olref[int][sem][pro]{prop:propositions-true} સાબિત કરો.
\end{prob}

\end{document}
